\documentclass[12pt,a4paper]{article}

\usepackage[utf8]{inputenc}
\usepackage[spanish,es-noshorthands]{babel}
\usepackage{geometry}
\geometry{left=2.5cm, right=2.5cm, top=3cm, bottom=2.5cm, headheight=15pt}
\usepackage{longtable}
\usepackage{booktabs}
\usepackage{array}
\usepackage{fancyhdr}
\usepackage{titlesec}
\usepackage{graphicx}
\usepackage{hyperref}
\usepackage{xcolor}
\usepackage{enumitem}
\usepackage{tabularx}
\usepackage{multirow}
\usepackage{makecell}
\usepackage{tikz}
\usetikzlibrary{arrows.meta, positioning, shapes.geometric}

\pagestyle{fancy}
\fancyhf{}
\fancyhead[L]{MANUAL DEL PROCESO}
\fancyhead[R]{ESTUDIO DE CASO EN INGENIERÍA DE SOFTWARE}
\fancyfoot[C]{\thepage}
\renewcommand{\headrulewidth}{0.4pt}
\renewcommand{\footrulewidth}{0.4pt}

\titleformat{\section}{\large\bfseries}{\thesection.}{0.5em}{}
\titleformat{\subsection}{\normalsize\bfseries}{\thesubsection}{0.5em}{}

\newcolumntype{L}[1]{>{\raggedright\arraybackslash}p{#1}}
\newcolumntype{C}[1]{>{\centering\arraybackslash}p{#1}}

\begin{document}

\begin{titlepage}
\centering
\vspace*{2cm}

{\Large\textbf{UNIVERSIDAD NACIONAL}}

\vspace{0.5cm}
{\Large\textbf{FACULTAD DE INGENIERÍA EN SISTEMAS COMPUTACIONALES}}

\vspace{3cm}

{\LARGE\textbf{MANUAL DEL PROCESO}}

\vspace{0.5cm}
{\LARGE\textbf{ESTUDIO DE CASO EN INGENIERÍA DE SOFTWARE}}

\vspace{2cm}

{\large\textbf{PROCESO DE INVESTIGACIÓN}}

\vspace{1cm}

\begin{tabular}{|l|l|}
\hline
\textbf{MACROPROCESO:} & INVESTIGACIÓN \\
\hline
\textbf{PROCESO:} & ESTUDIO DE CASO EN INGENIERÍA DE SOFTWARE \\
\hline
\textbf{SUBPROCESO:} & --- \\
\hline
\textbf{RESPONSABLE:} & INVESTIGADOR PRINCIPAL \\
\hline
\textbf{VERSIÓN:} & 1.0 \\
\hline
\end{tabular}

\vspace{3cm}

{\large Julio 2026}

\end{titlepage}

\tableofcontents
\newpage

\section{OBJETIVO}

Atender los requerimientos relacionados en base a una solicitud de estudio de caso en ingeniería de software, mediante la aplicación de un proceso estructurado que permita analizar, diseñar, ejecutar y documentar un caso de investigación que proporcione soluciones y conocimientos validados al dominio de la ingeniería de software.

\section{ALCANCE Y ÁMBITO DE APLICACIÓN}

Comprende desde la definición del objetivo de investigación, pasando por el diseño del protocolo, la recolección y análisis de datos, hasta la publicación de los resultados del estudio de caso.

El ámbito de aplicación del proceso es a nivel nacional, aplicable a proyectos de investigación en ingeniería de software que requieran un análisis profundo de casos específicos.

\section{NORMAS GENERALES DE OPERACIÓN}

\subsection{BASE LEGAL}

\begin{longtable}{|L{4cm}|C{1.5cm}|C{2cm}|L{6cm}|}
\hline
\textbf{NORMA} & \textbf{AÑO} & \textbf{ARTÍCULOS} & \textbf{TEMAS} \\
\hline
\endhead
Norma técnica de control interno emitida por la Contraloría General del Estado & 2014 & 410-08 & Adquisición de infraestructura tecnológica. \\
\hline
Resoluciones C.D. 457 & 2013 & Art. 3, num. 2.4.3 & Estructura Orgánica del IESS. (Dirección Nacional de Tecnología de la Información, funciones y responsabilidades) \\
\hline
Resolución CD.483 & 2015 & --- & Reformas al reglamento orgánico funcional del IESS, Resolución No. C.D.457 \\
\hline
\end{longtable}

\subsection{POLÍTICAS DEL PROCESO}

\begin{itemize}[noitemsep]
    \item Se deberá emitir un informe por cada estudio de caso realizado.
    \item Analizar el requerimiento conjuntamente con el área requirente.
    \item Los hallazgos deben ser validados antes de su publicación.
    \item Se debe mantener una cadena de evidencia que respalde los resultados.
    \item La participación de informantes requiere consentimiento informado.
\end{itemize}

\section{DESCRIPCIÓN DEL PROCESO}

\subsection{DIAGRAMA DE RELACIONAMIENTO DE LOS PROCESOS}

[No se incluye según parámetros establecidos]

\subsection{DIAGRAMA DE FLUJO}

El diagrama de flujo del proceso de Estudio de Caso en Ingeniería de Software se presenta a continuación, organizado en 5 fases según el modelo BPMN del Bizagi:

\begin{figure}[htbp]
\centering
\resizebox{0.95\textwidth}{!}{
\begin{tikzpicture}[node distance=1cm and 1cm, auto, font=\small, >={Stealth[length=3mm]}]
    % Estilos
    \tikzset{
        startstop/.style={rectangle, rounded corners, minimum width=2.5cm, minimum height=0.7cm, text centered, draw=black, fill=red!30},
        process/.style={rectangle, minimum width=2.2cm, minimum height=0.7cm, text centered, draw=black, fill=blue!20},
        decision/.style={diamond, minimum width=1.8cm, minimum height=0.7cm, text centered, draw=black, fill=green!20, aspect=2.5},
        arrow/.style={thick, ->},
        phase/.style={rectangle, minimum width=14cm, minimum height=0.5cm, text centered, draw=black, fill=gray!20, font=\bfseries}
    }
    
    % Fase 1: Diseño
    \node (fase1) [phase] {FASE 1: DISEÑO};
    \node (start) [startstop, below=0.3cm of fase1] {Inicio};
    \node (pro_pregunta) [process, below=of start] {Definir pregunta de investigación};
    \node (pro_objetivo) [process, below=of pro_pregunta] {Definir objetivo del estudio};
    \node (pro_seleccion) [process, below=of pro_objetivo] {Seleccionar caso(s) específico(s)};
    \node (gw_info) [decision, below=0.3cm of pro_seleccion] {¿Info suficiente?};
    
    % Fase 2: Preparación
    \node (fase2) [phase, below=1cm of gw_info] {FASE 2: PREPARACIÓN};
    \node (pro_metodos) [process, below=0.3cm of fase2] {Definir métodos de recolección};
    \node (pro_protocolo) [process, below=of pro_metodos] {Diseñar protocolo de investigación};
    \node (pro_procedimientos) [process, below=of pro_protocolo] {Definir procedimientos de campo};
    \node (pro_instrumentos) [process, below=of pro_procedimientos] {Preparar instrumentos de recolección};
    \node (pro_consentimiento) [process, below=of pro_instrumentos] {Obtener consentimiento informado};
    
    % Fase 3: Recolección
    \node (fase3) [phase, below=0.5cm of pro_consentimiento] {FASE 3: RECOLECCIÓN};
    \node (pro_piloto) [process, below=0.3cm of fase3] {Realizar estudio piloto};
    \node (gw_protocolo) [decision, below=0.3cm of pro_piloto] {¿Protocolo aprobado?};
    \node (pro_entrevistas) [process, below=0.3cm of gw_protocolo] {Realizar entrevistas semiestructuradas};
    \node (pro_observacion) [process, below=of pro_entrevistas] {Realizar observación directa con registros};
    \node (pro_artefactos) [process, below=of pro_observacion] {Recolectar artefactos y documentos};
    
    % Fase 4: Análisis
    \node (fase4) [phase, below=0.5cm of pro_artefactos] {FASE 4: ANÁLISIS};
    \node (pro_registrar) [process, below=0.3cm of fase4] {Registrar y organizar datos recolectados};
    \node (pro_codificar) [process, below=of pro_registrar] {Codificación de datos cualitativos};
    \node (pro_cadena) [process, below=of pro_codificar] {Construcción de cadena de evidencia};
    \node (pro_analisis) [process, below=of pro_cadena] {Análisis dentro/entre del caso(s)};
    \node (gw_hallazgos) [decision, below=0.3cm of pro_analisis] {¿Hallazgos validados?};
    
    % Fase 5: Reporte
    \node (fase5) [phase, below=1cm of gw_hallazgos] {FASE 5: REPORTE};
    \node (pro_informe) [process, below=0.3cm of fase5] {Redactar informe del caso};
    \node (pro_revision) [process, below=of pro_informe] {Revisión por pares académicos};
    \node (pro_retro) [process, below=of pro_revision] {Incorporar retroalimentación};
    \node (pro_publicar) [process, below=of pro_retro] {Publicar y diseminar resultados};
    \node (stop) [startstop, below=of pro_publicar] {Fin};
    
    % Flechas - Fase 1
    \draw [arrow] (start) -- (pro_pregunta);
    \draw [arrow] (pro_pregunta) -- (pro_objetivo);
    \draw [arrow] (pro_objetivo) -- (pro_seleccion);
    \draw [arrow] (pro_seleccion) -- (gw_info);
    
    % Flechas - Fase 1 a 2
    \draw [arrow] (gw_info.east) -- ++(1.5,0) node[anchor=south] {No} |- (pro_entrevistas.east);
    \draw [arrow] (gw_info) -- node[anchor=east] {Sí} (pro_metodos);
    
    % Flechas - Fase 2
    \draw [arrow] (pro_metodos) -- (pro_protocolo);
    \draw [arrow] (pro_protocolo) -- (pro_procedimientos);
    \draw [arrow] (pro_procedimientos) -- (pro_instrumentos);
    \draw [arrow] (pro_instrumentos) -- (pro_consentimiento);
    
    % Flechas - Fase 2 a 3
    \draw [arrow] (pro_consentimiento) -- (pro_piloto);
    \draw [arrow] (pro_piloto) -- (gw_protocolo);
    
    % Flechas - Gateway Protocolo
    \draw [arrow] (gw_protocolo) -- node[anchor=east] {Sí} (pro_entrevistas);
    \draw [arrow] (gw_protocolo.west) -- ++(-1.5,0) node[anchor=south] {No} |- (pro_instrumentos.west);
    
    % Flechas - Fase 3
    \draw [arrow] (pro_entrevistas) -- (pro_observacion);
    \draw [arrow] (pro_observacion) -- (pro_artefactos);
    
    % Flechas - Fase 3 a 4
    \draw [arrow] (pro_artefactos) -- (pro_registrar);
    
    % Flechas - Fase 4
    \draw [arrow] (pro_registrar) -- (pro_codificar);
    \draw [arrow] (pro_codificar) -- (pro_cadena);
    \draw [arrow] (pro_cadena) -- (pro_analisis);
    \draw [arrow] (pro_analisis) -- (gw_hallazgos);
    
    % Flechas - Gateway Hallazgos
    \draw [arrow] (gw_hallazgos) -- node[anchor=east] {Sí} (pro_informe);
    \draw [arrow] (gw_hallazgos.west) -- ++(-1.5,0) node[anchor=south] {No} |- (pro_entrevistas.west);
    
    % Flechas - Fase 5
    \draw [arrow] (pro_informe) -- (pro_revision);
    \draw [arrow] (pro_revision) -- (pro_retro);
    \draw [arrow] (pro_retro) -- (pro_publicar);
    \draw [arrow] (pro_publicar) -- (stop);
\end{tikzpicture}
}
\caption{Diagrama de Flujo del Proceso de Estudio de Caso en Ingeniería de Software (5 Fases)}
\end{figure}

\subsection{DESCRIPCIÓN DE ACTIVIDADES}

\begin{longtable}{|C{0.8cm}|L{3.5cm}|L{4cm}|C{2cm}|L{3.5cm}|}
\hline
\textbf{No.} & \textbf{ACTIVIDAD} & \textbf{DESCRIPCIÓN DE LA ACTIVIDAD} & \textbf{RESPONSABLE} & \textbf{DOCUMENTO DE REFERENCIA/SISTEMA} \\
\hline
\endhead

\multicolumn{5}{|c|}{\textbf{FASE 1: DISEÑO}} \\
\hline

1 & Definir pregunta de investigación & Se formula la pregunta o preguntas de investigación que guiarán el estableciendo el enfoque y alcance de la investigación. & Investigador Principal & Protocolo de Investigación \\
\hline

2 & Definir objetivo del estudio & Se establece claramente el objetivo del estudio de caso, definiendo qué se busca investigar y qué se espera lograr. & Investigador Principal & Protocolo de Investigación \\
\hline

3 & Seleccionar el/los caso(s) específico(s) & Se identifica y selecciona el caso o casos que serán estudiados, considerando criterios de relevancia y viabilidad. & Investigador Principal & --- \\
\hline

4 & ¿Información suficiente para el caso? & \multirow{2}{3.5cm}{Verificación de si se cuenta con información suficiente para continuar. Si es negativo, se vuelven a realizar entrevistas.} & --- & --- \\
\cline{1-4}
5 & Continuar recolección & & Equipo de Investigación & --- \\
\hline

\multicolumn{5}{|c|}{\textbf{FASE 2: PREPARACIÓN}} \\
\hline

6 & Definir métodos de recolección & Se determinan los métodos y técnicas que se utilizarán para la recolección de datos (entrevistas, observación, documentos). & Equipo de Investigación & --- \\
\hline

7 & Diseñar protocolo de investigación & Se elabora el documento protocolar que establece los procedimientos, cronograma y recursos necesarios. & Equipo de Investigación & Protocolo de Investigación \\
\hline

8 & Definir procedimientos de campo & Se establecen los procedimientos operativos para la recolección de datos en el lugar del estudio. & Equipo de Investigación & --- \\
\hline

9 & Preparar instrumentos de recolección & Se diseñan y preparan los instrumentos: guías de entrevista, formatos de observación, formularios. & Equipo de Investigación & Guía de Entrevista \\
\hline

10 & Obtener consentimiento informado & Se gestiona el consentimiento informado de todos los participantes del estudio, asegurando el cumplimiento ético. & Equipo de Investigación & Formulario de Consentimiento \\
\hline

\multicolumn{5}{|c|}{\textbf{FASE 3: RECOLECCIÓN}} \\
\hline

11 & Realizar estudio piloto & Se ejecuta una prueba piloto para validar los instrumentos y procedimientos antes de la recolección principal. & Equipo de Investigación & --- \\
\hline

12 & ¿Protocolo aprobado tras piloto? & \multirow{2}{3.5cm}{Verificación de si el protocolo fue validado satisfactoriamente durante el piloto. Si es negativo, se revisan los instrumentos.} & --- & --- \\
\cline{1-4}
13 & Revisar instrumentos & & Equipo de Investigación & --- \\
\hline

14 & Realizar entrevistas semiestructuradas & Se realizan entrevistas con los informantes clave del caso, siguiendo las guías de entrevista previamente diseñadas. & Equipo de Investigación & --- \\
\hline

15 & Realizar observación directa con registros & Se realiza observación directa del contexto y proceso del caso, manteniendo registros detallados. & Equipo de Investigación & --- \\
\hline

16 & Recolectar artefactos y documentos & Se recopilan todos los documentos y artefactos relevantes para el caso (manuales, diagramas, especificaciones). & Equipo de Investigación & Artefactos del caso \\
\hline

\multicolumn{5}{|c|}{\textbf{FASE 4: ANÁLISIS}} \\
\hline

17 & Registrar y organizar datos recolectados & Se sistematizan y organizan todos los datos obtenidos en la fase de recolección. & Equipo de Investigación & Registro de Auditoría (Audit Trail) \\
\hline

18 & Codificación de datos cualitativos & Se aplica un proceso de codificación a los datos cualitativos para facilitar su análisis. & Equipo de Investigación & Esquema de Codificación (Codebook) \\
\hline

19 & Construcción de cadena de evidencia & Se establece la relación entre hallazgos, evidencia y conclusiones para garantizar la trazabilidad. & Equipo de Investigación & Esquema de cadena \\
\hline

20 & Análisis dentro/entre del caso(s) & Se realiza el análisis de los datos tanto al interior de cada caso como entre los casos estudiados. & Investigador Principal & --- \\
\hline

21 & ¿Hallazgos validados satisfactoriamente? & \multirow{2}{3.5cm}{Verificación de la validez de los hallazgos obtenidos. Si es negativo, se recolectan más datos.} & --- & --- \\
\cline{1-4}
22 & Recolectar más datos & & Equipo de Investigación & --- \\
\hline

\multicolumn{5}{|c|}{\textbf{FASE 5: REPORTE}} \\
\hline

23 & Redactar informe del caso & Se elabora el informe final del estudio de caso incluyendo todos los hallazgos, análisis y conclusiones. & Investigador Principal & Informe Final del Estudio \\
\hline

24 & Revisión por pares académicos & El informe es revisado por expertos en el dominio para validar su contenido y calidad. & Investigador Principal & --- \\
\hline

25 & Incorporar retroalimentación & Se incorporan las observaciones y sugerencias de los revisores al documento. & Investigador Principal & --- \\
\hline

26 & Publicar y diseminar resultados & Se publican los resultados del estudio para su conocimiento y aplicación. & Investigador Principal & --- \\
\hline

\end{longtable}

\section{INDICADORES DE PROCESO}

\begin{longtable}{|C{0.8cm}|L{3cm}|L{4cm}|C{2.5cm}|C{2.5cm}|}
\hline
\textbf{No.} & \textbf{INDICADOR} & \textbf{DESCRIPCIÓN} & \textbf{META} & \textbf{UNIDAD DE MEDIDA} \\
\hline
\endhead

1 & Cumplimiento de plazos & Porcentaje de actividades ejecutadas dentro del plazo establecido & 90\% & Porcentaje \\
\hline

2 & Calidad del protocolo & Protocolo aprobado en primera revisión & 80\% & Porcentaje \\
\hline

3 & Satisfacción de informantes & Nivel de satisfacción de los participantes del estudio & 85\% & Porcentaje \\
\hline

4 & Calidad del informe final & Informe aprobado sin observaciones mayores & 90\% & Porcentaje \\
\hline

5 & Tiempo promedio de análisis & Tiempo promedio entre recolección y análisis & 15 días & Días \\
\hline

\end{longtable}

\section{GLOSARIO DE TÉRMINOS Y ABREVIATURAS}

\begin{longtable}{|L{4cm}|L{10cm}|}
\hline
\textbf{TÉRMINO/ABREVIATURA} & \textbf{DESCRIPCIÓN} \\
\hline
\endhead

Estudio de Caso & Investigación profunda de un caso específico en su contexto real, empleando múltiples fuentes de evidencia. \\
\hline

Protocolo de Investigación & Documento que establece los procedimientos, cronograma y recursos para el desarrollo del estudio. \\
\hline

Cadena de Evidencia & Documentación que permite seguir el rastro desde la pregunta inicial hasta las conclusiones finales. \\
\hline

Codificación & Proceso de etiquetar y categorizar datos cualitativos para facilitar su análisis. \\
\hline

Consentimiento Informado & Documento que certifica que los participantes conocen y aceptan voluntariamente participar en el estudio. \\
\hline

Triangulación & Uso de múltiples fuentes de datos o métodos para validar los hallazgos de la investigación. \\
\hline

Audit Trail & Registro detallado de todas las actividades realizadas durante el proceso de investigación. \\
\hline

Codebook & Esquema de codificación que define las categorías y sus códigos para el análisis de datos. \\
\hline

BPMN & Business Process Model and Notation - Notación de Modelado de Procesos de Negocio. \\
\hline

\end{longtable}

\section{ANEXOS}

\textbf{Anexo 1:} Protocolo de Investigación

\textbf{Anexo 2:} Guías de Entrevista

\textbf{Anexo 3:} Formulario de Consentimiento Informado

\textbf{Anexo 4:} Plantilla de Informe Final del Estudio

\textbf{Anexo 5:} Diagrama BPMN del Proceso (Archivo Bizagi: EstudioCaso.bpmn)

\end{document}
